\documentclass[10pt,a4paper]{amsart}
\usepackage[margin=19mm]{geometry}
\usepackage{amsmath,amssymb,mathtools}
\usepackage[scr=rsfs,cal=euler]{mathalfa}
\usepackage{xfrac}
\usepackage[unicode,hidelinks,bookmarks=false]{hyperref}
\newcommand{\ie}{i.e.\ }
\newcommand{\cf}{cf.\ }
\newcommand{\resp}{respectively\ }
\newcommand{\Subsection}{\subsection}
\providecommand{\nobreakdash}{\nobreak\hbox{-}}
\setlength{\emergencystretch}{2em}
\multlinegap0pt
\usepackage{amssymb}
\usepackage[shortlabels]{enumitem}
\usepackage{tikz}
\usepackage{xcolor}
\newtheorem{theorem}{Theorem}
\newtheorem{lemma}{Lemma}
\newcommand{\Q}{\mathbb{Q}}
\title{The Gompf $\theta$-Invariant of Canonical Contact Structures via Legendrian Surgery}
\author{Mohan Bhupal, Burak Ozbagci}
\date{Source: arXiv:2605.21152v1 (CC BY 4.0)}
\begin{document}
\maketitle

\noindent\emph{Source and licence.} The Gompf $\theta$-Invariant of Canonical Contact Structures via Legendrian Surgery. Version arXiv:2605.21152v1 (CC BY 4.0), distributed by arXiv under the Creative Commons Attribution 4.0 International licence, http://creativecommons.org/licenses/by/4.0/, as recorded in the arXiv licence metadata for this exact version and shown on the abstract page https://arxiv.org/abs/2605.21152v1. Full licence notice: \texttt{licenses/arxiv-2605.21152-CC-BY-4.0.txt}.\par\smallskip

\noindent\textit{Editorial scope.} Counted unit: the article's lemma lem:theta-sporadic-examples (source0.tex lines 2344--2354). The article's complete original proof of it is delivered with it (source0.tex lines 2356--2463, one \texttt{proof} environment of 108 lines). No mathematics is written, repaired or re-derived: the statement and the proof are the article's own lines, in the article's own notation, and no retained line is substituted or re-flowed.  Retained uncounted as the source-local context the proof needs: lines 787--808; lines 2247--2254.  Nothing was omitted from any retained span, so the package carries no omission record.  No retained line contains a citation, so the wrapper carries no bibliography and no bibliography entry.  No retained line contains a figure environment, an image-inclusion command or drawing code, so no image is delivered and the package declares no asset dependency.  Editorial formatting only, in addition: an AMS \texttt{amsart} preamble that restates the article's own declarations and macros used by the retained lines, namely the environments \texttt{theorem}, \texttt{lemma} on separate counters, together with the standard packages those lines need.  The retained environments print in this document as Theorem 1, Lemma 1 in order of appearance, whereas the article prints the same items with the numbering of its own full document.  The complete native source of the article as distributed in the arXiv e-print for this exact version is bundled unchanged as \texttt{sources/201043/source0.tex}.
\begin{theorem}\label{thm:main-seifert}
Let $Y=Y(e_0;r_1,\dots,r_k)$ be a Seifert fibered space with
normalized invariants $r_i\in\Q\cap(0,1)$, central weight
$e_0\le -1$, and negative rational Euler number
\[
e(Y)=e_0+r_1+\cdots+r_k<0.
\]
For each $i=1,\dots,k$, write
\[
\frac{1}{r_i}=\frac{p_i}{q_i}=[a_1^i,\dots,a_{n_i}^i],
\qquad a_j^i\ge 2,
\]
where $p_i,q_i$ are coprime integers with $0<q_i<p_i$. Then the
canonical contact structure $\xi_{\rm can}$ on $Y$ satisfies
\[
\theta(\xi_{\rm can})
=
(2k-1)
-\sum_{i=1}^k\!\left(I(p_i/q_i)+\frac{q_i+q_i^*+2}{p_i}\right)
+\frac{1}{e(Y)}\!\left(-e_0+k-2-\sum_{i=1}^k\frac{q_i+1}{p_i}\right)^{\!2}.
\]
\end{theorem}

\begin{lemma}\label{lem:theta-three-families}
Let $\ell \ge 1$, and let $\Gamma$ be a star-shaped plumbing graph with $k$ legs, each leg a linear chain of length $\ell$, all vertices having framing $-2$. Let the central vertex have framing $-b$. Then the canonical contact structure $\xi_{\rm can}$ on the Seifert fibered space $Y_\Gamma$ satisfies $\theta(\xi_{\rm can})=-2$ in each of the following cases:
\begin{enumerate}
\item[(1)] $k=4\ell+4$ and $b=k=4\ell+4$,
\item[(2)] $k=\ell^2+3\ell+3$ and $b=k+1=\ell^2+3\ell+4$,
\item[(3)] $k=4\ell^2-3$ and $b=k+4=4\ell^2+1$.
\end{enumerate}
\end{lemma}

\begin{lemma}\label{lem:theta-sporadic-examples}
Let $\Gamma$ be a star-shaped plumbing graph with $k$ legs, each leg a linear chain of length $\ell$, all vertices having framing $-2$, and let the central vertex have framing $-b$. Then the canonical contact structure $\xi_{\rm can}$ on the Seifert fibered space $Y_\Gamma$ satisfies $\theta(\xi_{\rm can})=-2$ in each of the following cases:
\begin{enumerate}
\item[(1)] $\ell=1$, $k=4$, $b=7$,
\item[(2)] $\ell=2$, $k=3$, $b=9$,
\item[(3)] $\ell=1$, $k=7$, $b=4$,
\item[(4)] $\ell=1$, $k=8$, $b=5$,
\item[(5)] $\ell=1$, $k=9$, $b=7$.
\end{enumerate}
Moreover, none of these examples is covered by Lemma~\ref{lem:theta-three-families}.
\end{lemma}

\begin{proof}
As in the proof of Lemma~\ref{lem:theta-three-families}, since each leg is a chain of length $\ell$ with all framings equal to $-2$, we have
\[
[2,\dots,2]=\frac{\ell+1}{\ell},
\]
so that
\[
p=\ell+1,\qquad q=\ell,\qquad q^*=\ell,\qquad I(p/q)=-\ell.
\]
Hence the general formula for $\theta(\xi_{\rm can})$ given in Theorem~\ref{thm:main-seifert} becomes
\[
\theta(\xi_{\rm can})
=
k\ell-1+\frac{(b-2)^2}{-b+\frac{k\ell}{\ell+1}}.
\]
Therefore $\theta(\xi_{\rm can})=-2$ is equivalent to
\begin{equation}\label{eq:sporadic-theta-condition}
(\ell+1)(b-2)^2=(k\ell+1)\big((\ell+1)b-k\ell\big).
\end{equation}
We now verify \eqref{eq:sporadic-theta-condition} in each case.

\medskip

\noindent
\textbf{Case (1):} $(\ell,k,b)=(1,4,7)$.

\[
(\ell+1)(b-2)^2=2\cdot 5^2=50,
\]
and
\[
(k\ell+1)\big((\ell+1)b-k\ell\big)
=(4+1)(2\cdot 7-4)=5\cdot 10=50.
\]
Thus \eqref{eq:sporadic-theta-condition} holds.

\medskip

\noindent
\textbf{Case (2):} $(\ell,k,b)=(2,3,9)$.

\[
(\ell+1)(b-2)^2=3\cdot 7^2=147,
\]
and
\[
(k\ell+1)\big((\ell+1)b-k\ell\big)
=(3\cdot 2+1)(3\cdot 9-3\cdot 2)=7\cdot 21=147.
\]
Thus \eqref{eq:sporadic-theta-condition} holds.

\medskip

\noindent
\textbf{Case (3):} $(\ell,k,b)=(1,7,4)$.

\[
(\ell+1)(b-2)^2=2\cdot 2^2=8,
\]
and
\[
(k\ell+1)\big((\ell+1)b-k\ell\big)
=(7+1)(2\cdot 4-7)=8\cdot 1=8.
\]
Thus \eqref{eq:sporadic-theta-condition} holds.

\medskip

\noindent
\textbf{Case (4):} $(\ell,k,b)=(1,8,5)$.

\[
(\ell+1)(b-2)^2=2\cdot 3^2=18,
\]
and
\[
(k\ell+1)\big((\ell+1)b-k\ell\big)
=(8+1)(2\cdot 5-8)=9\cdot 2=18.
\]
Thus \eqref{eq:sporadic-theta-condition} holds.

\medskip

\noindent
\textbf{Case (5):} $(\ell,k,b)=(1,9,7)$.

\[
(\ell+1)(b-2)^2=2\cdot 5^2=50,
\]
and
\[
(k\ell+1)\big((\ell+1)b-k\ell\big)
=(9+1)(2\cdot 7-9)=10\cdot 5=50.
\]
Thus \eqref{eq:sporadic-theta-condition} holds.

Therefore $\theta(\xi_{\rm can})=-2$ in all five cases.

Finally, none of these examples appears in Lemma~\ref{lem:theta-three-families}. Indeed, for $\ell=1$ that lemma gives only
\[
(k,b)=(8,8),\ (7,8),\ (1,5),
\]
and for $\ell=2$ it gives only
\[
(k,b)=(12,12),\ (13,14),\ (13,17).
\]
These do not include any of the five triples listed above.
\end{proof}

\end{document}
